\section{Methods}
\label{sec:methods}

To answer the question of ``does \openzl{} improve performance for data transfers
on HPC systems?'' we need to test various file transfer scenarios which can be
found in HPC workloads. Such scenarios include:
\begin{itemize}
\item large-size checkpointing
\item data ingest
\end{itemize}

Adding permutations to these scenarios, we can also vary when we compress the
data. For instance, we can compress the data:
\begin{itemize}
\item ``in-transit'' immediately before a transfer,
\item pre-compress before writing to storage, and
\item no compression.
\end{itemize}


However, it may not be practical, easy, or even possible to compress/decompress
in all these scenarios. For instance, data stored on network storage would need
to be compressed \emph{before} storage rather than ``in-transit'' since we don't
have control over the storage nodes. In which case, the only compression we can
test is pre-compressing data and storing that data.


\subsection{In-Transit}
\label{sec:in-transit-method}

\subsection{Checkpointing}
\label{sec:checkpoint-method}
When testing checkpointing compression, we gather our data, compress it, and
save it to a file. However, this also adds additional complexity/overhead
if/when we resume a job. Now we also have to \emph{decompress} the data when
ingesting it into the resuming job. Additionally, we need some sort of logic to
decide whether said data even needs to be decompressed. For instance, as
discussed in Section \ref{sec:compress-or-not}, not every file will be
compressed.

\subsection{No Compression}
\label{sec:no-compress-method}
Self-explanatory. We transmit the data without any compression whatsoever. If
our other compression scenarios are faster than this: great we have an
improvement.



\subsection{To Compress or not to Compress?}
\label{sec:compress-or-not}
Given the wide range of data formats, shapes, and sizes, even \emph{if}
\openzl{} is effective at compressing in a particular scenario, it may not be the
right choice in \emph{every} scenario. As such, we want to develop some method
of deciding whether or not compression will be beneficial \emph{without actually
doing the operation}.


\subsection{How Do We Quantify ``Effectiveness?''}
\label{sec:quantifying-effectiveness}
We've been asking the question: ``is \openzl{} effective for compression in HPC
workloads?'' But what does ``effective'' mean? Depending on what we're after,
whether or not it's ``effective'' could mean multiple things. Is it cheap to
implement? Does it save storage space? Does it save time? Does it work at all?
For our purposes, we mainly care about how much time it saves, if at all. As
such, we pretty much care about two metrics: how quickly can \openzl{} compress
our data? and what compression ratios can it achieve?

The faster the execution time, the better. Without compression, there is no
execution time, so \openzl{} is already at a disadvantage here. However, it may
save time in transmitting the data over the network by having fewer bytes to
transfer. This is where that second metric of compression ratios comes in.
